\documentclass[11pt,a4paper]{article}
\usepackage[utf8]{inputenc}
\usepackage[T1]{fontenc}
\usepackage[polish]{babel}
\usepackage{lmodern}
\usepackage[margin=2.3cm]{geometry}
\usepackage{booktabs}
\usepackage{array}
\usepackage{enumitem}
\usepackage{titlesec}
\usepackage{xcolor}
\usepackage{fancyhdr}
\usepackage{hyperref}

\definecolor{navy}{RGB}{27,42,86}
\definecolor{lgray}{RGB}{240,240,244}

\titleformat{\section}{\large\bfseries\color{navy}}{\thesection.}{0.5em}{}
\titleformat{\subsection}{\bfseries\color{navy}}{\thesubsection.}{0.5em}{}

\pagestyle{fancy}
\fancyhf{}
\renewcommand{\headrulewidth}{0.6pt}
\fancyhead[L]{\small Neurofizjologia i pomiar EEG}
\fancyhead[R]{\small Instrukcja laboratoryjna nr 4}
\fancyfoot[C]{\small\thepage}

\newcommand{\infobox}[1]{%
  \par\noindent\colorbox{lgray}{\parbox{\dimexpr\linewidth-2\fboxsep}{\vspace{2pt}#1\vspace{2pt}}}\par\vspace{4pt}}

\title{\vspace{-1.5cm}\textbf{\color{navy}Instrukcja laboratoryjna nr 4}\\[4pt]
\Large Metodologia badań psychofizjologicznych II: analiza potencjałów wywołanych i analiza częstotliwościowa}
\author{Neurofizjologia i pomiar EEG\\ Psychologia, studia II stopnia --- specjalność neuropsychologia kliniczna}
\date{Czas trwania zajęć: 120 minut}

\begin{document}
\maketitle
\thispagestyle{fancy}

\section{Cele zajęć}
\begin{itemize}[nosep]
  \item Wprowadzenie do metody potencjałów wywołanych (ERP) -- zasady segmentacji i uśredniania.
  \item Zapoznanie z głównymi komponentami ERP i ich interpretacją poznawczą.
  \item Wprowadzenie do analizy częstotliwościowej (widmowej) sygnału EEG -- transformata Fouriera i moc widmowa.
  \item Praktyczne wykonanie podstawowej analizy ERP oraz analizy widmowej na przykładowych danych.
\end{itemize}

\section{Wprowadzenie teoretyczne}

\subsection{Metoda potencjałów wywołanych (ERP)}
Pojedyncza odpowiedź mózgu na bodziec jest zwykle niewidoczna na tle spontanicznej aktywności EEG (stosunek sygnału do szumu jest niski). \textbf{Uśrednianie} wielu segmentów (\emph{epok}) zsynchronizowanych względem momentu prezentacji bodźca eliminuje aktywność niezwiązaną czasowo z bodźcem, uwydatniając stały w czasie komponent odpowiedzi -- potencjał wywołany.

\subsection{Główne komponenty ERP}
\begin{center}
\begin{tabular}{@{}lll@{}}
\toprule
\textbf{Komponent} & \textbf{Latencja} & \textbf{Interpretacja poznawcza}\\
\midrule
N100/P100 & 80--130 ms & Wczesne przetwarzanie sensoryczne, uwaga\\
N170 & ok. 170 ms & Przetwarzanie twarzy\\
N200 & 200--350 ms & Detekcja niezgodności, kontrola konfliktu\\
P300 (P3) & 250--500 ms & Ewaluacja bodźca, alokacja uwagi, aktualizacja pamięci\\
N400 & ok. 400 ms & Przetwarzanie semantyczne, niezgodność znaczeniowa\\
\bottomrule
\end{tabular}
\end{center}

\subsection{Analiza czasowo-częstotliwościowa i moc widmowa}
Podczas gdy ERP opisuje uśrednioną, zablokowaną fazowo odpowiedź na bodziec, \textbf{analiza widmowa} (np. transformata Fouriera, FFT) rozkłada sygnał ciągły na składowe częstotliwościowe, ujawniając moc poszczególnych pasm (delta--gamma) niezależnie od precyzyjnego zablokowania w fazie z bodźcem. \textbf{Analiza czasowo-częstotliwościowa} (np. transformata falkowa) łączy obie perspektywy, pokazując zmiany mocy widmowej w czasie względem zdarzenia (np. desynchronizację/synchronizację związaną ze zdarzeniem, ERD/ERS).

\subsection{Podstawy analizy źródłowej}
Analiza źródłowa (\emph{source localization}) próbuje odtworzyć, w oparciu o rozkład potencjału na powierzchni skóry głowy i model przewodnictwa tkanek, przybliżoną lokalizację generatorów korowych sygnału (np. metody dipolowe, sLORETA). Jest to zagadnienie odwrotne (\emph{inverse problem}) o niejednoznacznym rozwiązaniu, wymagające założeń upraszczających.

\infobox{\textbf{Do zapamiętania:} ERP wymaga wielu powtórzeń tego samego typu bodźca (zwykle $\geq$30--50 epok na warunek) ze względu na niski stosunek sygnału do szumu pojedynczej próby.}

\section{Sprzęt i oprogramowanie}
Komputery z zainstalowanym środowiskiem do analizy EEG (np. EEGLAB/MATLAB lub MNE-Python), przykładowy zbiór danych z paradygmatu oddball zarejestrowany na poprzednich zajęciach lub udostępniony przez prowadzącego.

\section{Przebieg zajęć (harmonogram, 120 min)}
\begin{center}
\begin{tabular}{@{}p{2.3cm}p{11.7cm}@{}}
\toprule
\textbf{Czas} & \textbf{Aktywność}\\
\midrule
0--20 min & Wykład: zasada uśredniania ERP, główne komponenty i ich interpretacja.\\
20--35 min & Wykład: podstawy analizy widmowej i czasowo-częstotliwościowej, wprowadzenie do analizy źródłowej.\\
35--55 min & Demonstracja: segmentacja danych, uśrednianie i wizualizacja ERP w oprogramowaniu.\\
55--65 min & Przerwa.\\
65--95 min & Ćwiczenie praktyczne: samodzielne wyznaczenie komponentu P300 na danych z zadania oddball (zadanie 1).\\
95--115 min & Ćwiczenie praktyczne: obliczenie mocy widmowej pasm dla warunku EO/EC (zadanie 2).\\
115--120 min & Dyskusja wyników, podsumowanie.\\
\bottomrule
\end{tabular}
\end{center}

\section{Zadania do wykonania}
\begin{enumerate}[label=\textbf{Zadanie \arabic*.}, leftmargin=2.7cm]
  \item Na udostępnionych danych z paradygmatu oddball wykonajcie segmentację epok, uśrednienie osobno dla bodźców standardowych i docelowych oraz zidentyfikujcie amplitudę i latencję komponentu P300 w elektrodzie Pz.
  \item Obliczcie i porównajcie moc widmową pasma alfa dla warunków EO i EC (efekt blokady alfa) na elektrodach potylicznych.
  \item Sporządźcie krótkie sprawozdanie (0,5--1 strony) opisujące uzyskane wyniki i ich zgodność z oczekiwaniami teoretycznymi.
\end{enumerate}

\section{Kryteria zaliczenia ćwiczenia}
Poprawność techniczna analizy ERP (35\%), poprawność analizy widmowej (35\%), jakość merytoryczna sprawozdania (30\%).

\section{Pytania kontrolne}
\begin{enumerate}[nosep]
  \item Dlaczego pojedyncza próba nie pozwala na wiarygodne zidentyfikowanie komponentu ERP?
  \item Jaką funkcję poznawczą odzwierciedla klasycznie komponent P300?
  \item Czym różni się analiza widmowa od analizy czasowo-częstotliwościowej?
  \item Na czym polega zagadnienie odwrotne w analizie źródłowej EEG?
\end{enumerate}

\section{Literatura}
\begin{enumerate}[nosep]
  \item Kalat, J. W. (2019). \emph{Biological Psychology}. Cengage Learning.
  \item Chaddad, A., Wu, Y., Kateb, R., \& Bouridane, A. (2023). Electroencephalography signal processing. \emph{Sensors, 23}(14), 6434.
  \item Biasiucci, A., Franceschiello, B., \& Murray, M. M. (2019). Electroencephalography. \emph{Current Biology, 29}(3), R80--R85.
\end{enumerate}

\end{document}
